\section{What would establish continued smooth existence?}
\label{sec:existence}

The question is whether prescribed smooth initial data generate a solution
of \eqref{eq:baseline} that continues to exist with controlled derivatives
as the propagation distance increases. Three claims must be separated:
existence on a specified finite interval; continuation for every finite
distance; and smoothness of all orders. A computation over a finite interval
can support the first claim empirically. Neither a bounded maximum amplitude
nor a bounded $L^2$ norm controls all derivatives.

\subsection{The available theory must match the equation}

Mosincat, Pilod and Saut prove local well-posedness in
$H^s(\mathbb R^2)$, $s>0$, and global well-posedness and scattering for
sufficiently small $L^2(\mathbb R^2)$ data for their \emph{temporal}, nonlocal
Dysthe equation; see their equation~(1.12) and Theorems~1.1--1.2
\cite{mps}. Their one-dimensional whole-line discussion concerns temporal
Dysthe with third-order dispersion and $H^s(\mathbb R)$, $s\geq1/4$.
These statements do not directly cover our second-order spatial equation
on a periodic profile domain. Nonzero constant backgrounds are also outside
the finite-$L^2(\mathbb R^2)$ hypothesis.

For one-dimensional periodic temporal Dysthe, Heller and Shao obtain
estimates and low-regularity ill-posedness results, but explicitly do not
prove local well-posedness of the full equation \cite[\S1.5]{heller}.
Belloni and Giuliani's recent global result for a local periodic
NLS--Dysthe family likewise does not settle our problem: their discussion
explicitly excludes the water-wave mean-flow term $u|D|(|u|^2)$
\cite[pp.~5--6]{belloni}. Thus we have not identified a theorem whose stated
hypotheses match \eqref{eq:baseline}. This is a limitation of the present
review, not a claim that all Dysthe existence questions are open.

\subsection{A rigorous finite-dimensional statement, and its limit}

Let $\TT$ be the periodic profile circle of length $L>0$, and let $P_K$
denote the orthogonal Fourier projection onto modes $|n|\leq K$.
Write $F(u)$ for the full right-hand side of \eqref{eq:baseline}.
The following statement concerns exact Galerkin evolution, with exact
nonlinear products and the stated Fourier mean-flow operator.

\begin{proposition}[Global existence at each fixed Fourier cutoff]
\label{prop:galerkin-global}
Let $K$ be a nonnegative integer, let $\epsilon\in\mathbb R$, and let
$u_K(0)$ belong to the range of $P_K$. The finite system
\[
  \partial_\xi u_K=P_KF(u_K)
\]
has a unique solution for every $\xi\in\mathbb R$, and
$\norm{u_K(\xi)}_{L^2(\TT)}=\norm{u_K(0)}_{L^2(\TT)}$.
\end{proposition}
\begin{proof}
Its Fourier coefficients solve a polynomial ordinary differential
equation in finitely many real variables, hence have a unique local
solution. Put $\rho=|u_K|^2$. Self-adjointness of $P_K$ and $P_Ku_K=u_K$
allow its removal from the mass derivative. The linear dispersive term
vanishes after integration by parts. The local cubic term and the mean-flow
term have zero real contribution because $\rho$ and $|D_\tau|\rho$ are real.
The remaining two derivative terms give
\[
 \frac{\dd}{\dd\xi}\norm{u_K}_{L^2}^2
 =2\operatorname{Re}\int_{\TT}\overline{u_K}F(u_K)\,\dd\tau
 =-10\epsilon\int_{\TT}\rho\,\partial_\tau\rho\,\dd\tau=0.
\]
Parseval's identity therefore bounds the coefficient vector throughout
its maximal interval. In finite dimension it remains in a compact ball,
where the polynomial vector field and its derivative are bounded.
The ordinary differential equation continuation theorem excludes a finite
endpoint in either direction.
\end{proof}

This proposition does not prove existence for the partial differential
equation, or stability of a chosen time-stepping method. For $s>0$ it gives
only
\[
 \norm{u_K(\xi)}_{H^s(\TT)}
 \leq\left(1+\left(\frac{2\pi K}{L}\right)^2\right)^{s/2}
       \norm{u_K(0)}_{L^2(\TT)},
\]
whose constant grows with $K$. Global finite systems can therefore fail
to provide the uniform regularity needed in the limit $K\to\infty$.
An aliased implementation need not equal this Galerkin system, and an
RK4 update does not inherit its exact mass conservation or global bound.

\subsection{The analytic continuation route}

A PDE proof must first establish local existence, uniqueness and
continuous dependence in a specified function space for the exact
equation and boundary conditions. Directly applying a finite-dimensional
ODE argument is insufficient: the nonlinear terms contain derivatives,
including the first-order nonlocal operator $|D_\tau|$, so a naive
same-$H^s$ nonlinear estimate loses derivatives. Suitable cancellations,
modified energy estimates, or dispersive estimates must recover this loss.

If a local theory supplies a lifespan bounded below in terms of the
$H^s$ norm, its restart argument gives a continuation criterion: a finite
maximal distance can occur only if that controlling norm becomes
unbounded. Establishing the needed norm bound uniformly on every finite
interval would then give global existence in that regularity class.
Passing from Galerkin approximations also requires sufficient compactness,
convergence of the derivative nonlinearities, and uniqueness of the limit.
For initially smooth data, a claim of smoothness of all orders additionally
needs persistence of higher regularity on the same interval; one fixed
$H^s$ estimate alone is not that claim.

Residual-based PINN error estimates also require PDE stability
\cite{mishra,bai}. For an exact solution $u$ and a candidate $v$, put
$E(\xi)=\norm{u(\xi)-v(\xi)}_{H^s}$ and
$r=\partial_\xi v-F(v)$. If a justified stability argument gives
$E'\leq CE+\norm{r}_{\mathcal R}$ almost everywhere, with $C\geq0$ and
an appropriate residual norm $\mathcal R$, then Gronwall's inequality gives
\[
 E(\xi)\leq e^{C\xi}\left(E(0)+
              \int_0^\xi\norm{r(\eta)}_{\mathcal R}\,\dd\eta\right).
\]
For this equation, derivative loss must actually be controlled and $C$
must be bounded on the relevant solution set, independently of a numerical
cutoff. Without that work, the displayed bound is conditional; an error
estimate presupposing an exact solution is not an existence proof.

\subsection{A rigorous route from a computed candidate}

A numerical solution or PINN can instead serve as the center of an
\emph{a posteriori} existence proof, following the general validated-numerics
approach \cite{validation}. On a fixed interval $[0,\Xi]$, choose
Banach spaces $X,Y$ and an initial-value defect map, for example
\[
 \mathcal G(v)=\bigl(\partial_\xi v-F(v),\ v(0)-u_0\bigr).
\]
The spaces must make this map well-defined and must encode periodicity
and the desired regularity. Complex-valued functions are treated as real
Banach spaces because the equation contains complex conjugation.

\begin{theorem}[Conditional contraction certificate]
\label{thm:certificate}
Let $U\subseteq X$ be open, let $\mathcal G:U\to Y$ be continuously
Fr\'echet differentiable, and let $A:Y\to X$ be bounded, real-linear and
injective. Suppose $\overline B_X(\widetilde u,r)\subseteq U$, $r>0$, and
verified bounds satisfy
\[
 \norm{A\mathcal G(\widetilde u)}_X\leq Y_0,\qquad
 \sup_{v\in\overline B_X(\widetilde u,r)}
 \norm{I_X-A D\mathcal G(v)}_{X\to X}\leq Z(r).
\]
If $Z(r)<1$ and $Y_0+Z(r)r\leq r$, then $\mathcal G$ has exactly one zero
in this closed ball.
\end{theorem}
\begin{proof}
Define $T(v)=v-A\mathcal G(v)$. The derivative bound and the mean-value
inequality show that $T$ is Lipschitz on the ball with constant $Z(r)$.
For $v$ in the ball,
$\norm{T(v)-\widetilde u}_X\leq Z(r)r+Y_0\leq r$.
The contraction theorem gives a unique fixed point. Injectivity of $A$
makes a fixed point equivalent to a zero of $\mathcal G$, proving both
existence and uniqueness within the ball.
\end{proof}

This is a sufficient criterion, not a certificate already obtained for
our equation. Constructing $X,Y,A$ and proving the displayed bounds is the
main mathematical work. The validation must cover all unresolved Fourier
modes, the nonlocal operator, roundoff with outward bounds, and the initial
condition including any initial tail. For successive slabs, it must carry
the certified endpoint uncertainty forward. A residual sampled on a grid,
or at PINN training points, supplies none of those uniform guarantees by
itself. Such a certificate establishes the specified finite interval;
global continuation still requires control over arbitrarily long intervals.
